\section{Improvement 2: a hybrid metric--parametric trial space}\label{sec:hybrid}
\subsection{Rationale and definition}
Distortion immunity requires $P_p(\bx)\subset S^e$; nothing else. The space $Q_p(\hat x,\hat
y)$ satisfies this requirement with $\tfrac12p(p+1)$ additional functions of total degree $p+1$
to $2p$ in physical coordinates. On a non-affine element the nodes $\bx_k$, images of a tensor
grid, do not form a tensor grid in $(\hat x,\hat y)$, and interpolation with these high-degree
tensor terms on the perturbed grid drives $\Lambda_S$ up; \cref{prop:structure} below makes this
precise, and \cref{fig:lebesgue} shows that removing these terms brings $\Lambda$ back to the GLL
level. The
parametric functions, in contrast, interpolate on the exact GLL grid in $\bxi$. The hybrid
space keeps metric functions only for $P_p$ and uses parametric functions for the remainder.

\begin{definition}[hybrid trial space]\label{def:hybrid}
With the Legendre polynomials $L_a$,
\begin{equation}
  C_p=\operatorname{span}\{L_a(\xi)L_b(\eta):\ 0\le a,b\le p,\ a+b>p\},\qquad
  \tilde T_p^e=\{\,q\circ\bx+r:\ q\in P_p(\hat x,\hat y),\ r\in C_p\,\}.
  \label{eq:hybrid}
\end{equation}
The trial functions are the $M_k\in\tilde T_p^e$ with $M_k(\bxi_l)=\delta_{kl}$, evaluated
through $\bxi$: $u_h(\bx(\bxi))=\sum_ku_kM_k(\bxi)$.
\end{definition}
$C_p$ is the $L^2(\hat\Omega)$-orthogonal complement of $P_p(\bxi)$ in $Q_p(\bxi)$, so
$Q_p(\bxi)=P_p(\bxi)\oplus C_p$; this fixes the complement canonically. (Another complement
$C'$ of $P_p(\bxi)$ in $Q_p(\bxi)$, e.g.\ $\operatorname{span}\{\xi^a\eta^b:a+b>p\}$, gives the same
space only if $C'\subset P_p(\bx)\circ\bx\oplus C_p$, which holds when $P_p(\bxi)\subset\tilde
T_p^e$, e.g.\ on bilinear elements, but not in general.) In the
implementation $P_p$ is represented by $L_a(\tilde s)L_b(\tilde t)$, $a+b\le p$, in the scaled
skew coordinates~\eqref{eq:bbox}, and the interpolation matrix is
\begin{equation}
  \bm V_T=\begin{bmatrix}\bm Q\\ \bm R\end{bmatrix},\qquad
  Q_{jk}=q_j(\bx_k),\quad R_{ik}=r_i(\bxi_k),
  \label{eq:VT}
\end{equation}
with the $d=\tfrac12(p+1)(p+2)$ basis functions $q_j$ of $P_p$ and the $n-d$ functions $r_i$
of $C_p$. Gradients: $\nabla_{\bx}(q\circ\bx)=\nabla_{\bx}q$ exactly, $\nabla_{\bx}r=\bm
J^{-\T}\nabla_{\bxi}r$. We denote the element with~\eqref{eq:Ke} and~\eqref{eq:hybrid} by \PH.

\subsection{Properties}
\begin{proposition}\label{prop:hybrid}
\begin{enumerate}
\item[(a)] $P_p(\bx)\circ\bx\subset\tilde T_p^e$ for every element map: \PH\ is complete and
  distortion immune for fields of total degree $\le p$, independently of the sign of $\det\bm J$.
\item[(b)] $\tilde T_p^e$ is frame invariant and contains the rigid-body modes.
\item[(c)] If $\bx(\bxi)$ is bilinear (straight-sided element) and $\bm V_T$ is non-singular,
  then $\tilde T_p^e=Q_p(\bxi)$, $M_k=N_k$, and \PH\ is the conventional (symmetric) spectral
  element. In particular, on affine elements \PH, \PM\ and \PP\ coincide.
\item[(d)] The trial functions are interpolatory, so the GLL nodal-quadrature mass matrix is
  diagonal.
\end{enumerate}
\end{proposition}
\begin{proof}
(a)~by definition. (b)~$P_p(\hat x,\hat y)=P_p(\bx)$ is invariant under rigid motions, and so is
$\bxi$; $P_1(\bx)$ contains the rigid-body modes. (c)~For bilinear $\bx$, $x,y\in Q_1(\bxi)$, so
a monomial of total degree $j\le p$ in $(x,y)$ is a product of $j$ functions in $Q_1(\bxi)$ and lies
in $Q_j(\bxi)\subset Q_p(\bxi)$ \cite{Arnold2002Approximation}; $C_p\subset Q_p(\bxi)$. Hence
$\tilde T_p^e\subset Q_p(\bxi)$, with equality of dimensions if $\bm V_T$ is non-singular.
(d)~$M_k(\bxi_l)=\delta_{kl}$.
\end{proof}
By (c), \PH\ differs from the conventional element only on elements with curved edges (or
non-bilinear interior node placement); it is unsymmetric only there. A consequence of the same
argument is that the conventional Lagrange spectral element already contains $P_p(\bx)$ on
straight-sided quadrilaterals, so its loss of accuracy under distortion stems from curved
(non-bilinear) element maps, such as those of the mid-side node distortions studied in
\cref{sec:numerics}.

\paragraph{Unisolvence.} Let $\Iref$ denote interpolation in $Q_p(\bxi)$ at the GLL grid, and
$\Ilow$ the projection onto the Legendre coefficients with $a+b\le p$.
\begin{proposition}[factorisation]\label{prop:factor}
Let $\bm V_{\mathrm{ref}}$ be the Legendre Vandermonde matrix of the GLL grid (rows
$L_a(\xi_k)L_b(\eta_k)$, all $a,b\le p$) and $\bm A\in\R^{d\times d}$ the matrix of the map
$q\mapsto\Ilow\Iref(q\circ\bx)$ on $P_p$, $A_{ij}$ = coefficient of the $i$-th low-order Legendre
product in $\Iref(q_j\circ\bx)$. Then
\begin{equation}
  |\det\bm V_T|=|\det\bm V_{\mathrm{ref}}|\;|\det\bm A| ,
  \label{eq:factor}
\end{equation}
so the nodes are unisolvent for $\tilde T_p^e$ iff $\bm A$ is non-singular.
\end{proposition}
\begin{proof}
Order the rows of $\bm V_{\mathrm{ref}}$ as low-order ($a+b\le p$) rows $\bm L$ followed by the
rows $\bm R$ of~\eqref{eq:VT}. Each row of $\bm Q$ is a vector of nodal values and has a unique
expansion in the rows of $\bm V_{\mathrm{ref}}$: $\bm Q=\bm A^{\T}\bm L+\bm B\bm R$, the
coefficients being those of $\Iref(q_j\circ\bx)$. Hence
$\bm V_T=\begin{bmatrix}\bm A^{\T}&\bm B\\ \bm0&\bm I\end{bmatrix}\begin{bmatrix}\bm L\\ \bm R
\end{bmatrix}$.
\end{proof}
On an affine element $q\circ\bx\in P_p(\bxi)$, $\Iref$ is exact, $\Ilow$ is the identity on
$P_p(\bxi)$, and $\bm A$ is the matrix of the invertible composition with an affine map. Since
$\det\bm A$ depends continuously on the nodes, the nodes are unisolvent for every map
sufficiently close to an affine one. For strongly non-affine maps unisolvence is not
guaranteed; \cref{sec:cylinder} shows that it fails on some folded elements at $p\ge7$, exactly
where $\bm A$ becomes singular. Since $\bm V_{\mathrm{ref}}$ is fixed and well conditioned,
$\mathrm{rcond}(\bm A)$ is a cheap element-level indicator ($d\times d$ instead of $n\times n$).

\paragraph{Structure of the interpolant and Lebesgue constant.}
\begin{proposition}\label{prop:structure}
If $\bm A$ is non-singular, the hybrid interpolant of nodal data $\bm f$ is
\begin{equation}
  \mathcal I_T\bm f=\Iref\bm f+(\mathrm{Id}-\Iref)(q_{\bm f}\circ\bx),\qquad
  q_{\bm f}=\sum_jc_jq_j,\quad \bm A\bm c=\Ilow\Iref\bm f ,
  \label{eq:structure}
\end{equation}
i.e.\ $M_k=N_k+(\mathrm{Id}-\Iref)(q_k\circ\bx)$ with $q_k=q_{\bm e_k}$. Consequently
\begin{equation}
  \Lambda_T\le\Lambda_{\mathrm{GLL}}+\|\bm E_{\bx}\|\,\|\bm A^{-1}\|\,\gamma_p ,
  \qquad \bm E_{\bx}:\bm c\mapsto(\mathrm{Id}-\Iref)\Bigl(\sum_jc_jq_j\circ\bx\Bigr),
  \label{eq:lebesgue-bound}
\end{equation}
with $\gamma_p=\|\Ilow\Iref\|_{\ell^\infty\to\ell^\infty}$ depending on $p$ only and the norms
$\ell^\infty\to L^\infty(\hat\Omega)$ for $\bm E_{\bx}$. On bilinear elements $\bm
E_{\bx}=\bm0$ and $\Lambda_T=\Lambda_{\mathrm{GLL}}$.
\end{proposition}
\begin{proof}
Write $\mathcal I_T\bm f=q\circ\bx+r$. Interpolation at the GLL grid is exact on $Q_p(\bxi)\ni
r$, so $\Iref\bm f=\Iref(q\circ\bx)+r$, i.e.\ $r=\Iref\bm f-\Iref(q\circ\bx)$, which lies in $C_p$
iff $\Ilow\Iref(q\circ\bx)=\Ilow\Iref\bm f$, i.e.\ $\bm A\bm c=\Ilow\Iref\bm f$. Substituting $r$ gives
\eqref{eq:structure}; \eqref{eq:lebesgue-bound} follows from the triangle inequality and
$\|\Iref\|=\Lambda_{\mathrm{GLL}}$. For bilinear maps $q\circ\bx\in Q_p(\bxi)$
(\cref{prop:hybrid}(c)), so $(\mathrm{Id}-\Iref)(q\circ\bx)=0$.
\end{proof}
\Cref{prop:structure} states that the hybrid element is the conventional spectral element
plus a correction that restores completeness, and that this correction is an
\emph{interpolation error} of polynomials composed with the element map. It is small whenever
the map is close to bilinear: for $\bx=\bx_{\mathrm{bil}}+\delta\,\bm g$ with a smooth
perturbation $\bm g$, $\|\bm E_{\bx}\|=O(\delta)$, since $q\circ\bx_{\mathrm{bil}}\in
Q_p(\bxi)$. For curved elements with smooth edges $q\circ\bx$ is a polynomial of degree $p$
times the degree of the map, and $(\mathrm{Id}-\Iref)(q\circ\bx)$ decays with the smoothness of
the map. The same argument applies to $Q_p(\hat x,\hat y)$, with $q\in Q_p(\hat x,\hat y)$, $\Ilow$
replaced by the identity and an $n\times n$ matrix $\bm A_Q$; but then the correction contains the
interpolation errors of $\hat x^a\hat y^b$ up to total degree $2p$, which do not vanish even for
bilinear maps ($\hat x^p\hat y^p\circ\bx\in Q_{2p}(\bxi)\setminus Q_p(\bxi)$ in general). This is
consistent with \cref{fig:lebesgue}: on the straight-sided trapezoid $\Lambda_T=\Lambda_{\mathrm{GLL}}$,
whereas $\Lambda_Q$ grows to $4\cdot10^7$ at $p=16$.

\paragraph{Element matrix.} In~\eqref{eq:Ke} the columns of $\bm B_R$ that belong to $P_p$
are polynomial in $\bxi$, so the Jacobian cancellation is retained for them; the columns of the
complement contain $\bm J^{-1}$, and the corresponding integrand is
$\propto(\operatorname{adj}\bm J^{\T}\nabla_{\bxi}N)^{\T}\bm D(\operatorname{adj}\bm
J^{\T}\nabla_{\bxi}r)/\det\bm J$ (in strain-matrix form), rational as for curved isoparametric elements. Polynomial fields are
unaffected: for $\bm u\in P_p(\bx)$ the interpolant is exact, $\sum_ku_k\nabla M_k=\nabla\bm u$
is polynomial, and the integrand $(\bm B_L^\ast)^{\T}\bm D\bm\varepsilon(\bm u)$ is the same as for
\PM. On folded elements the complement part has a singular integrand where $\det\bm J=0$ inside
the element; \PH\ is undefined if $\det\bm J=0$ at a quadrature point. The cost per element
equals that of \PM\ plus one evaluation of $\bm J^{-1}$ per point; the parametric coordinates of
the quadrature points are known in all element routines.
